\documentclass[11pt]{article}
\usepackage[margin=1in]{geometry}
\usepackage{enumitem}
% =============================================================================
% Preambles/header.tex — shared LaTeX/Beamer preamble for this project
%
% Use from a Beamer lecture by prepending:
%     \input{header}
% and compiling with TEXINPUTS=../Preambles:$TEXINPUTS (the /compile-latex
% skill does this automatically).
%
% PALETTE CONTRACT. Color names defined here MUST match the names in
% Quarto/theme-template.scss so Beamer and Quarto renderings use the same
% palette. The scripts/check-palette-sync.sh script enforces this.
%
% To customize: edit the HEX values below AND the matching $variable in
% Quarto/theme-template.scss, then run scripts/check-palette-sync.sh.
% =============================================================================

% -----------------------------------------------------------------------------
% Engine detection + encoding
%
% Our /compile-latex skill uses XeLaTeX, where inputenc is unnecessary and
% can produce encoding warnings. Only load inputenc under pdfTeX.
% -----------------------------------------------------------------------------
\usepackage{iftex}
\ifPDFTeX
  \usepackage[utf8]{inputenc}
\fi

\usepackage{amsmath, amssymb, amsthm}
\usepackage{booktabs}
\usepackage{graphicx}

% -----------------------------------------------------------------------------
% PALETTE — mirrors Quarto/theme-template.scss variable names.
% Load xcolor and define colors BEFORE anything that references them
% (notably hyperref, hypersetup, and the Beamer color assignments).
% -----------------------------------------------------------------------------
\usepackage{xcolor}

\definecolor{primary-blue}{HTML}{012169}      % headings, accents
\definecolor{primary-gold}{HTML}{B9975B}      % emphasis, borders
\definecolor{highlight-yellow}{HTML}{F2A900}  % markers, alerts
\definecolor{light-bg}{HTML}{E8EDF5}          % title / section backgrounds
\definecolor{jet}{HTML}{1A1A1A}               % body text

% Semantic colors (used in diagrams and annotations)
\definecolor{positive}{HTML}{15803D}          % good / observed / identified
\definecolor{negative}{HTML}{B91C1C}          % bad / problematic / bias
\definecolor{neutral}{HTML}{525252}           % reference / context

% Highlight accents (match .hi-* classes in the SCSS)
\definecolor{hi-slate}{HTML}{314F4F}
\definecolor{hi-green}{HTML}{2E8B57}
\definecolor{hi-red}{HTML}{C41E3A}

% Semantic aliases so skill templates and snippets can write
% \textcolor{accent}{...} without hard-coding "primary-blue".
\colorlet{accent}{primary-blue}
\colorlet{accent2}{primary-gold}

% -----------------------------------------------------------------------------
% Hyperref — loaded AFTER xcolor + palette so link colors resolve.
% -----------------------------------------------------------------------------
\usepackage{hyperref}
\hypersetup{colorlinks=true, linkcolor=primary-blue, urlcolor=primary-blue,
            citecolor=primary-blue, pdfborder={0 0 0}}

% -----------------------------------------------------------------------------
% Course code — set once per deck (after \input{header}, before \begin{document})
% via \coursecode{CS401}. Shown in the footer on every slide so a deck's course
% is identifiable without opening the title page. Empty by default.
% -----------------------------------------------------------------------------
\makeatletter
\newcommand{\coursecode}[1]{\gdef\@coursecodetext{#1}}
\gdef\@coursecodetext{}
\makeatother

% -----------------------------------------------------------------------------
% Beamer theme assignments (only apply under Beamer).
%
% We detect Beamer via \@ifclassloaded{beamer}, which is the documented,
% non-internal way. This must be wrapped in \makeatletter/\makeatother
% because \@ifclassloaded contains an @-catcode character.
% -----------------------------------------------------------------------------
\makeatletter
\@ifclassloaded{beamer}{%
  \usefonttheme{professionalfonts}
  \setbeamercolor{structure}{fg=primary-blue}
  \setbeamercolor{frametitle}{fg=primary-blue}
  \setbeamercolor{title}{fg=primary-blue}
  \setbeamercolor{subtitle}{fg=primary-gold}
  \setbeamercolor{author}{fg=primary-blue}
  \setbeamercolor{institute}{fg=neutral}
  \setbeamercolor{date}{fg=primary-gold}
  \setbeamercolor{itemize item}{fg=primary-blue}
  \setbeamercolor{itemize subitem}{fg=primary-gold}
  \setbeamercolor{alerted text}{fg=primary-gold}
  \setbeamercolor{block title}{fg=white, bg=primary-blue}
  \setbeamercolor{block body}{bg=light-bg}
  % Semantic block variants: block=definition/concept (blue), exampleblock=
  % worked example (green/positive), alertblock=key takeaway or caution (gold).
  % Keep to <=2 colored boxes per slide across ALL three variants (INV-7).
  \setbeamercolor{block title example}{fg=white, bg=positive}
  \setbeamercolor{block body example}{bg=light-bg}
  \setbeamercolor{block title alerted}{fg=white, bg=primary-gold}
  \setbeamercolor{block body alerted}{bg=light-bg}

  % Minimal footer: course code (left) + page X / N (right), no navigation clutter
  \setbeamertemplate{navigation symbols}{}
  \setbeamertemplate{footline}{%
    \color{neutral}\footnotesize\hspace{0.7em}\@coursecodetext\hfill
    \insertframenumber/\inserttotalframenumber\hspace{0.7em}\vspace{0.3em}}
}{}
\makeatother

% -----------------------------------------------------------------------------
% TikZ setup — libraries the snippet gallery uses
% -----------------------------------------------------------------------------
\usepackage{tikz}
\usetikzlibrary{arrows.meta, positioning, calc, decorations.pathreplacing,
                fit, shapes.geometric, backgrounds}

% Shared TikZ styles — snippets and hand-written diagrams can pick these up
% via `\begin{tikzpicture}[every node/.style={...}]` or by naming specific
% styles (e.g., `\node[dag-node] (X) at (0,0) {X};`).
\tikzset{
  % Node styles for causal DAGs and similar
  dag-node/.style = {
    draw=primary-blue, fill=light-bg, rounded corners=2pt,
    minimum width=1.2cm, minimum height=0.9cm, align=center, font=\small
  },
  decision-node/.style = {
    draw=primary-gold, fill=white, diamond, aspect=1.8,
    minimum width=1.8cm, minimum height=1.2cm, align=center, font=\small,
    inner sep=1pt
  },
  % Edge styles — observed vs counterfactual
  observed-edge/.style = {-{Latex[length=2.5mm]}, thick, draw=jet},
  counterfactual-edge/.style = {-{Latex[length=2.5mm]}, thick, draw=neutral, dashed},
  confound-edge/.style = {-{Latex[length=2.5mm]}, thick, draw=negative, dashed},
  % Data-point styles for plots
  observed-dot/.style = {fill=primary-blue, draw=primary-blue, circle, inner sep=1.2pt},
  counterfactual-dot/.style = {fill=white, draw=neutral, circle, inner sep=1.2pt}
}

% -----------------------------------------------------------------------------
% Convenience macros
% -----------------------------------------------------------------------------
\newcommand{\muted}[1]{\textcolor{neutral}{#1}}
\newcommand{\key}[1]{\textcolor{primary-gold}{\textbf{#1}}}
\newcommand{\good}[1]{\textcolor{positive}{#1}}
\newcommand{\bad}[1]{\textcolor{negative}{#1}}

% Standout frame macro: full-bleed dark background for section transitions.
% Beamer-only (uses \begin{frame}/\setbeamercolor/\usebeamerfont) -- do not
% call from an article-class file (e.g. Notes/<CODE>/*-notes.tex). \muted,
% \key, \good, \bad, and \coursecode above are plain xcolor/gdef and are
% safe to use from either class; verified when Notes/article-class support
% was added.
% Expand: \transitionslide{Your transition title}
\newcommand{\transitionslide}[1]{%
  \begingroup
  \setbeamercolor{background canvas}{bg=primary-blue}
  \begin{frame}[plain]
    \centering
    \vfill
    {\usebeamerfont{title}\color{white}\Large #1\par}
    \vfill
  \end{frame}
  \endgroup
}

% Section divider frame (Beamer-only), an alternative to \transitionslide with a
% darker two-line style: near-black (jet) background, a small white label line
% above a larger gold title, and a thin gold rule along the bottom edge.
% Expand: \sectiondivider{Section 2}{Pointers}
\newcommand{\sectiondivider}[2]{%
  \begingroup
  \setbeamercolor{background canvas}{bg=jet}
  \begin{frame}[plain]
    \vspace*{3.0cm}
    {\color{white}\ttfamily\large #1\par}
    \vspace{0.5cm}
    {\color{primary-gold}\LARGE #2\par}
    \begin{tikzpicture}[remember picture, overlay]
      \draw[primary-gold, line width=2.5pt]
        ($(current page.south west)+(0,0.1)$) -- ($(current page.south east)+(0,0.1)$);
    \end{tikzpicture}
  \end{frame}
  \endgroup
}

% -----------------------------------------------------------------------------
% Channel watermark — "Concepts That Click" (YouTube).
%
% Article-class files (CompetitiveExam Q/A sets, Notes, Assignments) get a
% light diagonal draft watermark on every page plus a centered footer naming
% the channel. Beamer decks get a subtle TikZ background watermark instead
% (the footline already carries the course code). Centralized here so every
% MCQ PDF is branded without per-file edits.
% -----------------------------------------------------------------------------
\makeatletter
\@ifclassloaded{article}{%
  \usepackage{draftwatermark}
  \SetWatermarkText{Concepts That Click}
  \SetWatermarkScale{1}
  \SetWatermarkFontSize{2cm}
  \SetWatermarkLightness{0.86}
  \SetWatermarkAngle{45}
  \usepackage{fancyhdr}
  \pagestyle{fancy}
  \fancyhf{}
  \fancyfoot[C]{\footnotesize\textcolor{neutral}{Concepts That Click~|~YouTube: \href{https://www.youtube.com/@concepts.thatclick}{@concepts.thatclick}}}
  \renewcommand{\headrulewidth}{0pt}
  \renewcommand{\footrulewidth}{0pt}
  \fancypagestyle{plain}{%
    \fancyhf{}%
    \fancyfoot[C]{\footnotesize\textcolor{neutral}{Concepts That Click~|~YouTube: \href{https://www.youtube.com/@concepts.thatclick}{@concepts.thatclick}}}%
    \renewcommand{\headrulewidth}{0pt}%
    \renewcommand{\footrulewidth}{0pt}%
  }
}{}
\@ifclassloaded{beamer}{%
  \setbeamertemplate{background}{%
    \begin{tikzpicture}[remember picture, overlay]
      \node[rotate=45, opacity=0.055, align=center]
        at (current page.center)
        {\Huge\textcolor{neutral}{Concepts That Click}};
    \end{tikzpicture}%
  }
}{}
\makeatother 
\coursecode{JKSSB}
\renewcommand{\thesection}{04.\arabic{section}}
\newtheorem{example}{Example}[section]
\newenvironment{definitionbox}[1]{%
  \par\noindent\textbf{Definition (#1).}\ \itshape
}{\par\vspace{0.3em}}
\title{Organization of a Computer --- Detailed Lecture Notes}
\author{JKSSB Computer Awareness}
\date{\today}

\begin{document}
\maketitle

\noindent\textbf{Video lesson:}
\href{https://youtu.be/92XQ_n1t_cc}{How Computers Actually Work: CPU, Memory \& I/O Explained}

\section{The functional units of a computer}
Start with the office desktop. A clerk types marks on the keyboard, the CPU
processes them through a spreadsheet formula, the monitor and printer present
the result, and the file is kept on the SSD for later use. This single machine
contains every functional unit that the exam asks you to name.

\begin{definitionbox}{Functional unit}
A major part of a computer that performs one broad job --- accepting input,
producing output, holding data, performing operations, or controlling the
machine.
\end{definitionbox}

\begin{center}
\begin{tabular}{p{0.25\textwidth}p{0.63\textwidth}}
\toprule
\textbf{Unit} & \textbf{Job in this lesson}\\
\midrule
Input unit & Accepts data and instructions from the outside world.\\
Output unit & Presents the result to the user.\\
Memory (primary storage) & Holds data and instructions currently in use.\\
Arithmetic Logic Unit (ALU) & Performs arithmetic and logical operations.\\
Control unit & Directs the sequence of operations and the flow of data.\\
\bottomrule
\end{tabular}
\end{center}

The memory unit described here is \emph{primary storage}, that is, main
memory. Secondary storage such as a hard disk or an SSD is a separate idea and
is treated in Section 4.

\section{The CPU: ALU, control unit and registers}
The \key{Central Processing Unit (CPU)} is the processing centre of the
computer. It is built from three parts: the arithmetic logic unit, the control
unit, and registers.

\begin{definitionbox}{Central Processing Unit (CPU)}
The unit that fetches, decodes and executes instructions; it is formed by the
ALU, the control unit and the registers.
\end{definitionbox}

\begin{definitionbox}{Arithmetic Logic Unit (ALU)}
The part of the CPU that performs arithmetic operations (such as addition and
subtraction) and logical operations (such as comparison) on binary data.
\end{definitionbox}

\begin{definitionbox}{Control unit}
The part of the CPU that fetches and decodes instructions, issues control
signals, and coordinates the operation of all other units.
\end{definitionbox}

\begin{definitionbox}{Register}
A small, very fast storage location inside the CPU used to hold data,
addresses or control information during processing.
\end{definitionbox}

The ALU does the calculating, but it does not decide what to do. The control
unit decides the sequence of operations and the flow of data, and then tells
the ALU and the other units when to act. Registers are the CPU's own working
storage; they are not the same as main memory or secondary storage.

\begin{example}
In the marks spreadsheet, the formula that adds three marks is carried out by
the ALU. The control unit has already fetched the formula's instruction from
memory, decoded it, and signalled the ALU to act. The intermediate values are
kept in registers while the operation runs. The final total is then written
back to memory and later shown on the monitor.
\end{example}

\section{Input and output units}
The input unit brings data and instructions into the system. The output unit
presents the result.

\begin{center}
\begin{tabular}{p{0.30\textwidth}p{0.58\textwidth}}
\toprule
\textbf{Input devices} & \textbf{Output devices}\\
\midrule
Keyboard, mouse, scanner, microphone, light pen, digitizer & Monitor,
printer, plotter, speaker, sound-card output\\
\bottomrule
\end{tabular}
\end{center}

A device is classified by the role it plays in the described task, not by its
shape. A monitor is normally an output device because it presents results. A
scanner is an input device because it brings data in. When a question lists
devices, mark each one as input or output before choosing the option.

\section{Memory (primary storage) versus secondary storage}
This is a high-value contrast.

\begin{definitionbox}{Primary storage (main memory)}
The memory the CPU accesses directly while processing; it holds the programs
and data currently in use. Random Access Memory (RAM) is volatile; Read Only
Memory (ROM) is non-volatile.
\end{definitionbox}

\begin{definitionbox}{Secondary storage}
Storage that keeps data and files permanently but is not accessed directly by
the CPU during processing; examples are the hard disk, the Solid State Drive
(SSD) and the pen drive.
\end{definitionbox}

The key rule for the exam is that the \key{CPU works with main memory, not
with secondary storage}. When a program on the hard disk is run, it is first
loaded into main memory; the CPU then works from main memory. Secondary
storage is reached through main memory and the operating system.

\begin{example}
The spreadsheet file is saved on the SSD. When the clerk opens it, the file is
first loaded into RAM. All the additions and comparisons are then performed
on data held in RAM and in the CPU registers. Writing the result back to the
SSD happens later, through the operating system. At no point does the ALU
process the file straight from the SSD.
\end{example}

\section{The system bus: address, data and control}
The \key{system bus} is the set of wires that carries information between the
CPU, memory and the other units. It has three parts, and their directions
differ --- this is a common trap.

\begin{center}
\begin{tabular}{p{0.20\textwidth}p{0.26\textwidth}p{0.46\textwidth}}
\toprule
\textbf{Bus} & \textbf{Direction} & \textbf{Role}\\
\midrule
Address bus & CPU $\rightarrow$ memory & Carries the address of the memory
location the CPU wants to use.\\
Data bus & Bidirectional & Carries the actual data and instructions between
the CPU and memory.\\
Control bus & CPU $\rightarrow$ units & Carries control and timing signals,
such as read and write.\\
\bottomrule
\end{tabular}
\end{center}

\begin{definitionbox}{System bus}
The address bus, data bus and control bus taken together; they connect the
CPU with memory and the input/output units.
\end{definitionbox}

Remember the three questions: the address bus answers \emph{where}, the data
bus answers \emph{what}, and the control bus answers \emph{when and how}. The
address bus is one-way (from the CPU outwards) and carries the address, not
the contents of the location. The data bus is two-way. The control bus
carries signals, not data or addresses.

\begin{example}
The CPU wants the value stored at memory location 200. It places the address
200 on the address bus and sends a read signal on the control bus. The memory
puts the contents of location 200 onto the data bus, and the CPU receives it.
The address travelled out on the address bus; the data travelled back on the
data bus.
\end{example}

\section{The motherboard, expansion slots and interfaces}
The \key{motherboard} is the main circuit board that connects and links the
CPU, main memory and the buses. Expansion slots on the motherboard allow
extra cards to be added, such as a graphics card or a network card. Interfaces
and ports connect external devices --- keyboard, printer, monitor --- to the
system. The system bus runs along the motherboard between the CPU and the
other units, so the motherboard is the physical platform on which the units
communicate.

\section{The stored-program (von Neumann) idea}
In early machines, changing the task often meant rewiring the machine. The
\key{stored-program concept} changed this.

\begin{definitionbox}{Stored-program concept (von Neumann architecture)}
The design in which a program's instructions are stored in main memory
together with the data, so the CPU can fetch and execute them without
rewiring the machine.
\end{definitionbox}

Because instructions and data share the same memory, the CPU fetches an
instruction from memory, executes it, and moves on to the next. This repeating
pattern is the instruction cycle, described next.

\section{The instruction cycle}
The \key{instruction cycle} is the sequence the CPU follows for each
instruction.

\begin{enumerate}
  \item \key{Fetch}: the CPU obtains the next instruction from main memory.
  \item \key{Decode}: the control unit works out what the instruction means.
  \item \key{Execute}: the ALU or another unit performs the operation.
  \item \key{Store} (when required): the result is written back to memory or
    to a register.
\end{enumerate}

The order is fixed: fetch, then decode, then execute. The optional store step
writes the result back. Data and instructions move between the units over the
system bus, and the whole cycle repeats for the next instruction.

\begin{example}
Consider an instruction ``add the value in memory to a register''. During
fetch, the CPU gets the instruction from memory. During decode, the control
unit identifies it as an addition and signals the ALU. During execute, the ALU
performs the addition. The updated value is then stored back in the register
or in memory, and the cycle begins again.
\end{example}

Three distinctions prevent most errors in this topic:
\begin{enumerate}
  \item The CPU accesses main memory directly, not secondary storage.
  \item Address, data and control buses are not interchangeable: the address
    bus is one-way and carries addresses; the data bus is two-way and carries
    data and instructions; the control bus carries control signals.
  \item Memory (primary) holds current work; storage (secondary) keeps files
    permanently.
\end{enumerate}

\subsection*{Exercises}
\begin{enumerate}[label=\arabic*.]
  \item Name the five functional units of a computer and state one job of
    each.
  \item What are the three parts of the CPU, and what does each do?
  \item State the direction and role of the address bus, the data bus and the
    control bus.
  \item A file is opened from a hard disk and edited. Explain why the CPU does
    not process the file directly from the hard disk.
  \item List the steps of the instruction cycle in order.
\end{enumerate}

\subsection*{Solutions}
\begingroup\small
\begin{enumerate}[label=\arabic*.]
  \item Input unit (accepts data and instructions); output unit (presents
    results); memory/primary storage (holds data and instructions in use);
    ALU (performs arithmetic and logical operations); control unit (directs
    operations and data flow).
  \item The ALU, the control unit and the registers. The ALU computes; the
    control unit fetches, decodes and signals; the registers hold data and
    addresses during processing.
  \item Address bus: one-way from CPU to memory, carrying the address of the
    location. Data bus: bidirectional, carrying data and instructions.
    Control bus: carries control and timing signals from the CPU to other
    units.
  \item The CPU accesses only main memory directly. The file is first loaded
    into main memory, the CPU processes the copy in main memory and registers,
    and the operating system writes the changes back to the hard disk.
  \item Fetch, decode, execute (and store when a result must be written back).
\end{enumerate}
\endgroup
\end{document}
